\section{概率计算}

\answer{PG1-1}
记 $E$ 为至少出现一次正面，$C=i$ 为硬币类型。给定类型，$E$ 的补事件是连续 $n$ 次反面，因此似然为 $L_i=1-(1-p_i)^n$。全概率与贝叶斯公式给出
\[
\boxed{w:=\Prob(C=1\mid E)=\frac{\pi L_1}{\pi L_1+(1-\pi)L_2}.}
\]
给定类型后下一次抛掷与已知事件独立，所以
\[
\boxed{\Prob(H_{n+1}\mid E)=wp_1+(1-w)p_2.}
\]
当 $n>1$ 时，记录合并了许多不同序列，其似然是 $L_i$，通常不等于单次正面的似然 $p_i$。随着 $n\to\infty$，两种似然都趋于 $1$，后验趋回先验；因为两种硬币最终几乎必然都至少出现一次正面。

\answer{PG1-2}
记 $s=\alpha+\beta$，$E$ 为已知事件。$A$ 在 $u$ 附近响而 $B$ 尚未响的联合密度为 $\alpha e^{-su}$。因此
\[
\Prob(E)=\int_0^t\alpha e^{-su}\dd u=\frac\alpha s(1-e^{-st}).
\]
相除得
\[
\boxed{f_{T\mid E}(u)=\frac{s e^{-su}}{1-e^{-st}},\quad0<u<t.}
\]
分部积分计算条件均值：
\[
\E(T\mid E)=\frac{-t e^{-st}+(1-e^{-st})/s}{1-e^{-st}}
=\boxed{\frac1s-\frac{t}{e^{st}-1}}.
\]
当 $t\to\infty$ 时，密度趋于参数为 $s$ 的指数密度，均值趋于 $1/s$。已知胜出时钟本身不改变首次响铃时间，但附加的期限条件将其截断。

\answer{PG1-3}
记 $d=1-a-b$。给定 $T=1$ 与 $T=0$，去掉两者共同的多项式系数后，似然分别为 $a^hb^zd^e$ 与 $b^ha^zd^e$。因此
\[
\boxed{O:=\frac{\Prob(T=1\mid\text{记录})}{\Prob(T=0\mid\text{记录})}
=\frac\pi{1-\pi}\left(\frac ab\right)^{h-z}.}
\]
后验概率为 $O/(1+O)$。缺失因子 $d^e$ 抵消，所以在本模型中，同样的 $h,z$ 对应同样后验，与缺失个数无关。

若两类真值下缺失概率分别为 $d_1,d_0$，每个缺失记录会贡献似然比 $d_1/d_0$，不再自动抵消；非缺失输出的概率也须使用相应真值下的实际概率重新计算。因而不能只保留原公式而忽略缺失机制。

\answer{PG1-4}
令 $S=(X+Z)/\sqrt2$、$D=(X-Z)/\sqrt2$、$b=a/\sqrt2$。$S,D$ 独立标准正态，筛选只限制 $|S|<b$，其概率为 $2\Phi(b)-1$。由对称性条件均值为零，分部积分给出
\[
v:=\Var(S\mid |S|<b)
=\frac{\int_{-b}^b s^2\phi(s)\dd s}{2\Phi(b)-1}
=1-\frac{2b\phi(b)}{2\Phi(b)-1}.
\]
$D$ 的方差仍为 $1$，且与截断后的 $S$ 独立。因此
\[
\boxed{\operatorname{Corr}(X,Z\mid |X+Z|<a)=\frac{v-1}{v+1}.}
\]
当 $b\downarrow0$，条件下 $|S|<b$ 直接给出 $0\le v\le b^2$，故相关趋于 $-1$；当 $b\to\infty$，截断消失，$v\to1$，相关趋于零。

\answer{PG2-1}
第一次不可能重复，且第四次必然重复，故 $T\in\{2,3,4\}$。前两次不同的概率为
\[
\Prob(T>2)=1-\left[(1/2)^2+2(1/4)^2\right]=\frac58.
\]
前三次不同要求三种类型各出现一次，共有 $3!$ 种顺序，因此
\[
\Prob(T>3)=6\cdot\frac12\cdot\frac14\cdot\frac14=\frac3{16}.
\]
由 $\Prob(T>0)=\Prob(T>1)=1$，尾和公式给出
\[
\boxed{\E T=1+1+\frac58+\frac3{16}=\frac{45}{16},
\qquad \Prob(T=4)=\frac3{16}.}
\]
最后一个等式使用 $T\le4$；它不是某一指定类型出现的概率。

\answer{PG2-2}
每个红段恰有一个从黑色到红色的边界。对任意一条按固定方向读取的相邻关系，其起点黑、终点红的概率为 $nm/[N(N-1)]$。由 $N$ 个边界指示变量求和，
\[
\boxed{\E(\text{红段数})=\frac{mn}{N-1}.}
\]
若全部红卡连续，只需指定唯一红段的起点，就能确定整条颜色序列，共 $N$ 种。因两种颜色均存在，这些起点产生不同的标号序列。总颜色序列数为 $\binom Nm$，所以
\[
\boxed{\Prob(\text{只有一个红段})=\frac N{\binom Nm}.}
\]
$m=1$ 或 $n=1$ 时该概率等于 $1$，与直接观察一致。

\answer{PG2-3}
对每个三元子集 $A$，令 $I_A$ 表示它同桶。则 $\E I_A=s_3$，所以 $\E C=\binom m3s_3$。不同三元组的协方差按交集大小分类：不相交时独立；共享两个对象时，两个事件同时发生要求四个对象同桶，协方差为 $s_4-s_3^2$；共享一个对象时要求五个对象同桶，协方差为 $s_5-s_3^2$。

共享两个对象的无序三元组对有 $6\binom m4$ 个：在四个对象中选两个为公共部分。共享一个对象的有 $15\binom m5$ 个：先选公共对象，再把其余四个分成两对，共 $5\cdot3=15$ 种。方差展开中每个无序对出现两次，故
\[
\boxed{\begin{aligned}
\E C&=\binom m3s_3,\\
\Var C&=\binom m3s_3(1-s_3)
+12\binom m4(s_4-s_3^2)\\
&\quad+30\binom m5(s_5-s_3^2).
\end{aligned}}
\]
约定当 $r>m$ 时 $\binom mr=0$，所以公式也适用于 $m=3,4$。

\answer{PG2-4}
令 $Y_i$ 在普通卡 $i$ 位于全部目标之前时为 $1$，位于全部目标之后时为 $-1$，其余为零。反转排列给出 $\E Y_i=0$。在该卡与 $K$ 张目标卡的相对次序中，首尾各有概率 $1/(K+1)$，故 $\E Y_i^2=2/(K+1)$。

对 $i\ne j$，两卡同在目标之前的概率为 $2/[(K+1)(K+2)]$，同在之后也相同；指定 $i$ 在前、$j$ 在后的概率为 $1/[(K+1)(K+2)]$，反向亦同。因此
\[
\E(Y_iY_j)=\frac2{(K+1)(K+2)}.
\]
写 $D=\sum_iw_iY_i$，展开平方并用 $2\sum_{i<j}w_iw_j=(\sum_iw_i)^2-\sum_iw_i^2$，得到
\[
\boxed{\E D=0,\qquad
\Var D=\frac{2[(K+1)\sum_iw_i^2+(\sum_iw_i)^2]}{(K+1)(K+2)}.}
\]
这是代数恒等式，不要求权重非负；$m=0$ 时两边均为零。

\answer{PG3-1}
先由单纯形体积求等待次数的尾概率：
\[
\Prob(N>k)=\Prob(U_1+\cdots+U_k\le1)=\frac1{k!},\quad k\ge0.
\]
所以 $\E N=\sum_{k\ge0}1/k!=e$。把成本写为非负项之和，交换求和与期望，
\[
\E C=\sum_{j\ge1}\E[U_j^2\ind\{N\ge j\}].
\]
事件 $\{N\ge j\}$ 等价于前 $j-1$ 个增量的和不超过 $1$，只依赖过去，因此与新的 $U_j$ 独立。每一项等于 $\E U_j^2\Prob(N\ge j)=\Prob(N\ge j)/3$。于是
\[
\boxed{\E C=\frac13\E N=\frac e3.}
\]
$N$ 确实依赖整列增量；这里使用的是“第 $j$ 次抽样是否发生”与该次尚未观察的增量独立，而不是把整个停止次数当作独立变量。

\answer{PG3-2}
记 $t=-\log x$、$b=\log2$。负对数长度为 $nb+G_n$，其中 $G_n=\sum_{i=1}^n(-\log U_i)$，其密度为 $s^{n-1}e^{-s}/(n-1)!$。由于 $G_n>0$ 几乎必然，只要 $nb\ge t$ 就已越界，因此 $N\le L:=\lceil t/b\rceil$ 几乎必然。

对 $n\ge1$，若 $a_n=t-nb\le0$，则 $\Prob(N>n)=0$；若 $a_n>0$，
\[
\boxed{\Prob(N>n)=\Prob(G_n\le a_n)
=1-e^{-a_n}\sum_{j=0}^{n-1}\frac{a_n^j}{j!}.}
\]
最后一步由伽马密度反复分部积分得到。再加上 $\Prob(N>0)=1$，可写
\[
\boxed{\E N=1+\sum_{n=1}^{L-1}
\left[1-e^{-(t-nb)}\sum_{j=0}^{n-1}\frac{(t-nb)^j}{j!}\right].}
\]
$x=1/4$ 时 $L=2$。第一步越界等价于 $U_1/2<1/4$，概率为 $1/2$；否则第二步必越界。因此 $\Prob(N=1)=\Prob(N=2)=1/2$，$\E N=3/2$。

\answer{PG3-3}
对 $0<s<1$，所有正累计和的密度之和为 $u(s)=\sum_{k\ge1}s^{k-1}/(k-1)!=e^s$。首次进入 $(a,b]$ 的位置为 $s$ 时，前一步可从 $S_0=0$ 出发，也可从某个正累计和 $t\le a$ 出发。因 $s-t\in(0,1)$，新增量密度为 $1$，所以首次进入的密度为
\[
1+\int_0^a e^t\dd t=e^a,\quad a<s\le b.
\]
记 $d=b-a$。若恰有 $k\ge1$ 次进入，所有这些累计和依次为 $a<s_1<\cdots<s_k\le b$。内部各转移密度均为 $1$，最后一步越过 $b$ 的概率为 $1-b+s_k$。先固定 $s_k=a+v$，其前面的 $k-1$ 个位置形成体积 $v^{k-1}/(k-1)!$ 的单纯形，因此
\[
\begin{aligned}
\Prob(K=k)
&=\frac{e^a}{(k-1)!}\int_0^d v^{k-1}(1-d+v)\dd v\\
&=\boxed{e^a\left(\frac{d^k}{k!}-\frac{d^{k+1}}{(k+1)!}\right)}.
\end{aligned}
\]
对 $k\ge1$ 求和望远镜相消为 $e^ad$，故 $\boxed{\Prob(K=0)=1-e^ad}$。它非负，因为 $d\le1-a$ 且 $e^a(1-a)\le1$。最后由计数期望等于密度积分，
\[
\boxed{\E K=\int_a^bu(s)\dd s=e^b-e^a.}
\]

\answer{PG3-4}
未来缩小比例与已知历史独立。按题 \ref{prob:P24} 的归一化表示，给定当前历史 $\mathcal F_n$，
\[
Q=Q_n+(1-R_n)D',
\]
其中 $D'$ 独立于 $\mathcal F_n$，且 $\E D'=1$、$\Var D'=1/2$。因此
\[
\boxed{\widehat Q=\E(Q\mid\mathcal F_n)=Q_n+1-R_n,}
\quad
\boxed{\Var(Q\mid\mathcal F_n)=\frac{(1-R_n)^2}{2}.}
\]
最优性由条件均值的正交性质证明：对任意基于历史且二阶可积的预测 $Z$，
\[
\E(Q-Z)^2=\E\Var(Q\mid\mathcal F_n)+\E(\widehat Q-Z)^2,
\]
交叉项因条件期望为零而消失。故取 $Z=\widehat Q$ 达到最小值。又 $1-R_n=(1-x)\prod_{i=1}^nU_i$，所以
\[
\boxed{\E(Q-\widehat Q)^2=\frac{(1-x)^2}{2\cdot3^n}.}
\]
部分和 $Q_n$ 漏掉了条件期望为 $1-R_n$ 的正尾项；补上它才能消除可预测的偏差。

\answer{PG4-1}
引入 $S=U+V$，则 $U=WS$、$V=(1-W)S$，雅可比为 $s$。单位正方形给出的取值域为
\[
0<w<1,\qquad 0<s<\frac1{\max(w,1-w)}.
\]
因此
\[
\boxed{f_W(w)=\int_0^{1/\max(w,1-w)}s\dd s
=\frac1{2\max(w,1-w)^2},\quad0<w<1.}
\]
密度不是常数，所以不均匀。交换 $U,V$ 将 $W$ 变为 $1-W$，故 $\E W=1/2$。二阶矩分段积分：
\[
\begin{aligned}
\E W^2&=\frac12\int_0^{1/2}\frac{w^2}{(1-w)^2}\dd w
+\frac12\int_{1/2}^1 1\dd w\\
&=\left(\frac34-\log2\right)+\frac14=1-\log2.
\end{aligned}
\]
所以 $\boxed{\Var W=3/4-\log2}$。份额对称只决定均值，不决定其密度恒定。

\answer{PG4-2}
在 $0<l<h<1$ 上，极值的联合密度为 $n(n-1)(h-l)^{n-2}$：选定最小与最大样本后，其余 $n-2$ 个必须落在二者之间。令 $r=h-l$、$c=(l+h)/2$，逆变换雅可比为 $1$，且 $0<r<2\min(c,1-c)$。因此
\[
\boxed{f_C(c)=\int_0^{2\min(c,1-c)}n(n-1)r^{n-2}\dd r
=n[2\min(c,1-c)]^{n-1},\quad0<c<1.}
\]
密度关于 $1/2$ 对称，故 $\E C=1/2$。利用对称性并令 $u=2c$，
\[
\begin{aligned}
\Var C&=2\int_0^{1/2}(1/2-c)^2n(2c)^{n-1}\dd c\\
&=\frac n4\int_0^1(1-u)^2u^{n-1}\dd u
=\boxed{\frac1{2(n+1)(n+2)}}.
\end{aligned}
\]
最后的积分可展开 $(1-u)^2$，得到 $1/n-2/(n+1)+1/(n+2)$。

\answer{PG4-3}
记间距为 $D_1,\ldots,D_n$。给定全部间距，随机点落在第 $i$ 段的概率是 $D_i$，而不是 $1/n$。对非负函数 $g$，
\[
\E g(L)=\E\sum_iD_i g(D_i)=n\E[D_1g(D_1)].
\]
单个间距满足 $\Prob(D_1>l)=(1-l)^{n-1}$，密度为 $(n-1)(1-l)^{n-2}$，故
\[
\boxed{f_L(l)=n(n-1)l(1-l)^{n-2},\quad0<l<1.}
\]
这是 $\operatorname{Beta}(2,n-1)$ 密度。由贝塔积分，
\[
\boxed{\E L=\frac2{n+1},\qquad
\E(1/L)=n(n-1)\int_0^1(1-l)^{n-2}\dd l=n.}
\]
等概率选段的平均长度为 $\E(\sum_iD_i/n)=1/n$，严格小于 $2/(n+1)$。区别来自长段被随机点选中的概率更大。

\answer{PG4-4}
向量方向在圆周上均匀。两个方向相差 $\theta$ 时，投影异号对应总角度为 $2\theta$ 的两个扇形，因此 $\Prob(D_1=1)=\Prob(D_2=1)=\theta/\pi$。

设三个实际投影值为 $V_0,V_1,V_2$。三角恒等式给出
\[
V_1=\frac{V_0+V_2}{2\cos\theta}.
\]
因 $\cos\theta>0$，若外侧两个投影同号，中间投影也同号。故不可能同时发生两次符号变化，即 $\Prob(D_1=D_2=1)=0$。零投影为零概率事件。结合边际概率，
\[
\boxed{\begin{array}{c|cc}
& D_2=0&D_2=1\\ \hline
D_1=0&1-2\theta/\pi&\theta/\pi\\
D_1=1&\theta/\pi&0
\end{array}}
\]
于是
\[
\boxed{\Cov(D_1,D_2)=-\frac{\theta^2}{\pi^2},\qquad
\Prob(A=B=C)=1-\frac{2\theta}{\pi}.}
\]
相邻比较共享中间投影，两次符号变化因此不是独立事件。

\par\Needspace{8\baselineskip}
\subsection*{小节练习}
\answer{OP02}
把乙前 $2n$ 次的正面数记为 $Z$，最后一次的正面指标记为 $B$，于是 $Y=Z+B$，其中 $B$ 独立且以相同概率取 $0,1$。令 $D=Z-2X$。将甲和乙前 $2n$ 次的每个正反面全部互换，样本的概率不变，而
\[
D\longmapsto(2n-Z)-2(n-X)=-D.
\]
故 $\Prob(D>0)=\Prob(D<0)$。当 $D>0$ 时必然获胜；当 $D<0$ 时，即使 $B=1$ 也不可能使 $D+B>0$；当 $D=0$ 时仅在 $B=1$ 时获胜。因此
\[
\Prob(Y>2X)=\Prob(D>0)+\frac12\Prob(D=0)
=\boxed{\frac12}.
\]
对称性作用于整个差值 $Z-2X$，并不要求 $Z$ 与 $2X$ 同分布。

\answer{OP30}
记 $F$ 为选中双正面硬币，$H$ 为两次均为正面。先验概率和似然分别为
\[
\Prob(F)=\Prob(F^c)=\frac12,\qquad
\Prob(H\mid F)=1,\quad \Prob(H\mid F^c)=\frac14.
\]
由全概率公式，$\Prob(H)=\frac12+\frac12\cdot\frac14=\frac58$。因此
\[
\Prob(F\mid H)=\frac{\Prob(F)\Prob(H\mid F)}{\Prob(H)}
=\frac{1/2}{5/8}=\boxed{\frac45}.
\]
两次观测都来自同一枚已选定的硬币，故不能每次重新混合两类硬币的概率。

\answer{OP18}
把尚未闭合的连通部分看作开放链。初始有 $n$ 条开放链，每条恰有两个自由端。若当前有 $k$ 条开放链，则共有 $2k$ 个自由端：一对端点若属于同一条链，就闭合一个环；否则把两条链合成一条。两种情形都使开放链数减少 $1$，故过程依次且恰好经历 $k=n,n-1,\ldots,1$。

在 $k$ 条开放链的状态下，所有 $\binom{2k}{2}$ 对端点中，恰有 $k$ 对属于同一条链。因此该步产生新环的条件概率总是
\[
\frac{k}{\binom{2k}{2}}=\frac1{2k-1},
\]
与此前形成了哪些链无关。令 $I_k$ 表示这一步是否闭环，则最终环数 $C=\sum_{k=1}^n I_k$。利用期望线性性，
\[
\boxed{\E C=\sum_{k=1}^n\frac1{2k-1}.}
\]
求期望无需假设各步闭环事件独立。

\answer{OP25}
任意两条弦的四个端点几乎必然互异。条件于四个端点的圆周顺序，它们形成两条弦的三种无序配对等可能。只有连接交错端点的那一种配对在圆内部相交，因此每对弦相交的概率为 $1/3$。

记 $I_{ij}$ 为第 $i,j$ 条弦在圆内部相交的指标。连续取样下，端点重合、弦重合以及三条弦共点均是零概率事件；例如先固定两条相交弦，再取第三条弦的一个端点，其另一个端点只有至多一个位置能使它通过既有交点。故交点数几乎必然等于
\[
K=\sum_{1\le i<j\le n}I_{ij}.
\]
无需假设这些指标独立，直接得到
\[
\boxed{\E K=\binom n2\frac13=\frac{n(n-1)}6.}
\]
此处的随机弦模型由两个独立均匀端点定义；换一种选弦方式可能改变答案。

\answer{OP26}
反过来数最后一张 A 后面剩下的普通牌。对每张非 A 牌 $i$，令 $I_i$ 表示它是否排在全部四张 A 之后。这张牌与四张 A 的相对次序均匀；在这五张牌中，它成为最后一张的概率为 $1/5$，故 $\E I_i=1/5$。

共有 $48$ 张非 A 牌，因此最后一张 A 后的牌数 $R$ 满足
\[
R=\sum_{i=1}^{48}I_i,\qquad \E R=\frac{48}{5}.
\]
又因为 $T+R=52$，
\[
\boxed{\E T=52-\frac{48}{5}=\frac{212}{5}=42.4.}
\]
这里计算的是包括第四张 A 的停止位置，而不是第四张 A 之前的牌数；后者还要减去 $1$。

\answer{OP22}
在模式尚未失效、当前值为 $x$ 时，记 $f(x)$ 为下一次必须上升时还需生成的期望次数，$g(x)$ 为下一次必须下降时的相应期望。下一次总要生成一个数；若比较成功，才继续。因此
\[
f(x)=1+\int_x^1g(y)\dd y,\qquad
 g(x)=1+\int_0^xf(y)\dd y.
\]
这些期望有限：若前 $n$ 个数仍满足模式，则每对 $(U_{2k-1},U_{2k})$ 都必须上升，故其概率至多为 $2^{-\lfloor n/2\rfloor}$。同样的几何尾界控制从任一状态开始的剩余次数，足以支持上述积分方程。

微分得 $f'=-g$、$g'=f$，于是 $f''+f=0$。由 $f(1)=1$、$g(0)=1$，即 $f'(0)=-1$，有
\[
f(x)=A\cos x-\sin x,\qquad
A=\frac{1+\sin1}{\cos1}=\sec1+\tan1.
\]
第一次生成之后需要上升，故
\[
\begin{aligned}
\E T&=1+\int_0^1 f(x)\dd x\\
&=A\sin1+\cos1
=\boxed{\sec1+\tan1}.
\end{aligned}
\]
三角函数的角度单位为弧度。

\answer{OP24}
先从任意起点 $R_0=x\in[0,1)$ 考虑同一过程。可写为
\[
1-R_k=(1-x)\prod_{j=1}^kU_j,
\qquad U_j\overset{\mathrm{iid}}\sim\operatorname{Unif}(0,1).
\]
有限乘积 $Q_n=\prod_{k=1}^nR_k$ 递减且在 $[0,1]$ 内，故收敛到 $Q$。令 $g(x)=\E[Q\mid R_0=x]$，并约定 $g(1)=1$。由有界收敛，题目所求为 $g(0)$。

为了正确确定 $g$ 在右端点的边界，利用 $0\le1-\prod a_k\le\sum(1-a_k)$ 得
\[
0\le1-g(x)\le\sum_{k=1}^{\infty}\E(1-R_k)
=(1-x)\sum_{k=1}^{\infty}2^{-k}=1-x.
\]
所以 $g(x)\to1$ 当 $x\uparrow1$，不能把 $g$ 在 $[0,1)$ 上恒取零。

按第一次抽样条件化，并用有界收敛处理后续无限乘积，
\[
(1-x)g(x)=\int_x^1t g(t)\dd t.
\]
右边是有界函数的积分，先推出 $g$ 在 $[0,1)$ 连续，再可求导，得到
\[
-g(x)+(1-x)g'(x)=-xg(x),\qquad g'(x)=g(x).
\]
因此 $g(x)=Ce^x$。由已证明的 $g(1-)=1$ 得 $C=e^{-1}$，最终
\[
\boxed{\lim_{n\to\infty}\E\prod_{k=1}^nR_k=e^{-1}.}
\]

\answer{OP35}
给定最大值 $H=h$，其余 $n-1$ 个样本在排序前相当于从 $(0,h)$ 独立均匀抽样。故对 $0\le t\le h$，条件最小值满足
\[
\Prob(L>t\mid H=h)=\left(1-\frac th\right)^{n-1},
\]
从而
\[
\E[L\mid H=h]=\int_0^h\left(1-\frac th\right)^{n-1}\dd t
=\frac hn.
\]
因此 $\E L=\E H/n$，且
\[
\E[LH]=\E\!\left[H\E(L\mid H)\right]=\frac1n\E H^2,
\qquad \Cov(L,H)=\frac1n\Var(H).
\]
把全部样本变为 $1-X_i$ 不改变联合分布，却把最小值变为 $1-H$，所以 $L\overset d=1-H$，从而 $\Var(L)=\Var(H)>0$。最终
\[
\boxed{\Corr(L,H)=\frac1n.}
\]
条件分布均按连续密度所确定的版本理解；不把零概率事件 $H=h$ 当作普通事件比值。

\answer{OP36}
记各段长度为 $D_1,\ldots,D_n$，其和为 $1$。正长度线段能够组成非退化多边形，当且仅当最大长度严格小于其余长度之和，因此本题条件是
\[
\max_iD_i<\frac12.
\]
必要性来自沿其余各边连接最大边两端时的严格三角不等式。充分性还可保证构型无自交。记最大边长为 $L$，总长为 $S$，假设 $2L<S$。对半径 $r\ge L/2$ 的圆，长度 $D_i$ 对应的较小圆心角为
\[
\theta_i(r)=2\arcsin\frac{D_i}{2r}.
\]
若在 $r=L/2$ 时这些角之和至少为 $2\pi$，随半径增大，角和连续降至零，故存在半径使角和恰为 $2\pi$。否则，将最大边对应的角改用优弧角 $2\pi-\theta_L(r)$：需要其余角之和等于 $\theta_L(r)$。两者之差在 $r=L/2$ 时为负，而充分大的 $r$ 下与 $(S-2L)/r$ 同号，因而为正。由连续性仍存在所需半径。两种情形都可依次在圆周上截取这些正角，再连接相邻顶点，得到边长合要求的凸多边形，因此无自交。

排序后的均匀切点联合密度为常数 $(n-1)!$。变换为段长后，密度仍在单纯形 $D_i>0$、$\sum D_i=1$ 上为常数，故各段长交换对称。第一段大于 $1/2$ 当且仅当所有切点都在 $1/2$ 右边，所以
\[
\Prob(D_i>1/2)=\Prob(D_1>1/2)=2^{-(n-1)}.
\]
总长为 $1$，不可能同时有两段大于 $1/2$；任一段恰等于 $1/2$ 的概率为零。于是
\[
\boxed{\Prob(\text{能够组成多边形})=1-\frac{n}{2^{n-1}}.}
\]
特别地，$n=3$ 时概率为 $1/4$，$n=4$ 时为 $1/2$。

\answer{OP37}
将点排序为 $X_{(1)}<\cdots<X_{(N)}$，令内部间距
\[
D_i=X_{(i+1)}-X_{(i)},\qquad 1\le i\le N-1.
\]
互为最近邻的两点必定相邻。对每一条相邻边，令 $I_i$ 表示两端互为最近邻。端点边 $i=1$ 只需满足 $D_1<D_2$，末端边同理；内部边 $2\le i\le N-2$ 则需满足
\[
D_i<D_{i-1},\qquad D_i<D_{i+1}.
\]
均匀次序统计量的所有间距（连同两端边界间距）交换对称，且相等概率为零。因此端点边成为配对的概率为 $1/2$，内部边的三个间距中自身最小的概率为 $1/3$。

每个点只有一个最近邻，所以不同互选配对没有公共端点；每个配对恰贡献两个好点。若 $G$ 为好点数，则
\[
G=2\sum_{i=1}^{N-1}I_i,
\quad
\E G=2\left(\frac12+\frac12+\frac{N-3}{3}\right)
=\boxed{\frac{2N}{3}}.
\]
$N=3$ 时内部边求和为空，上式仍成立。这里数的是点；互选配对数的期望为 $N/3$。

